% Main file for Thematic Paper plan, Traditional Format.
% PR07284 allows flexible chapter structure at advisor/committee discretion.
\documentclass{muthesis2026}
\thematicplan
\usepackage[numbers]{natbib}
\setcounter{secnumdepth}{4}
\usepackage{graphicx}
\graphicspath{{figures/}}
\usepackage{latexsym}
\usepackage{amsmath}
\usepackage{amsthm}
\usepackage{url}
\usepackage{textcomp}
\usepackage{caption,setspace}
\captionsetup{font={stretch=1.5}}
\usepackage[flushleft]{threeparttable}

\makeatletter
\renewcommand{\@dotsep}{10000}
\setlength{\@fptop}{0pt}
\setlength{\@fpbot}{0pt plus 1fil}
\makeatother

\title{TITLE OF THEMATIC PAPER}
\author{STUDENT FULL NAME}
\input{preamble}

\begin{document}

\maketitle

\linespacing{1.5}
\acknowledgements{Write acknowledgements here in formal written language. Length must not exceed one page.}

\abstract{%
\linespacing{1.2}
Write the abstract here. Keep within two pages.
}
{
\linespacing{1.2}
Write the implication of the thematic paper here.
}

\tableofcontents
\listoftables
\listoffigures

% The chapter structure below is a suggestion. Adjust as appropriate
% to the discipline and nature of the research, per advisor discretion.

\chapter{Introduction}
Describe background, significance, objectives, and scope \cite{mount2004bioinformatics}.

\chapter{Literature Review}
Review relevant concepts, theories, and related studies.

\chapter{Research Methodology}
Describe research procedures and methods.

\chapter{Results and Discussion}
Present and discuss findings.

\chapter{Conclusion and Recommendations}
Summarize key findings, limitations, and future work.

\bibliographystyle{unsrtnat}
\bibliography{references}

\biography

\end{document}
